\documentclass[11pt,oneside]{article}

\pdfminorversion=7
\usepackage{QuizStyle}

\hypersetup{
    pdftitle={20260929 Quiz for MATH301 Modern Algebra I},
    pdfauthor={Jungwoo Kim},
    pdfsubject={Quiz for MATH301 Modern Algebra I},
    pdfcreator={LaTeX}
}

\begin{document}

\quiztitle{20260929 Quiz}

\begin{exercise}{1.7.1}
Let $F$ be a field. Show that the multiplicative group of nonzero elements of $F$, denoted by $F^\times$, acts on the set $F$ by
\[
    g\cdot a=ga,
\]
where $g\in F^\times$, $a\in F$, and $ga$ is the usual product in $F$ of the two field elements. State clearly which axioms in the definition of a field are used.
\end{exercise}

\begin{solution}
For $g_1,g_2\in F^\times$ and $a\in F$, associativity of multiplication gives
\[
    g_1\cdot(g_2\cdot a)=g_1(g_2a)=(g_1g_2)a=(g_1g_2)\cdot a.
\]
The multiplicative identity gives
\[
    1\cdot a=a.
\]
Hence $F^\times$ acts on $F$. The action axioms use associativity of multiplication and the multiplicative identity. The field axioms also ensure that the nonzero elements are closed under multiplication and inverses, so $F^\times$ is a group.
\end{solution}

\begin{exercise}{1.7.16}
Let $G$ be any group and let $A=G$. Show that the maps defined by
\[
    g\cdot a=gag^{-1}
\]
for all $g,a\in G$ satisfy the axioms of a left group action. This action of $G$ on itself is called \emph{conjugation}.
\end{exercise}

\begin{solution}
For $g_1,g_2,a\in G$,
\[
\begin{aligned}
    g_1\cdot(g_2\cdot a)
    &=g_1(g_2ag_2^{-1})g_1^{-1}\\
    &=(g_1g_2)a(g_1g_2)^{-1}\\
    &=(g_1g_2)\cdot a.
\end{aligned}
\]
Also
\[
    1\cdot a=1a1^{-1}=a.
\]
Thus conjugation defines a left group action of $G$ on itself.
\end{solution}

\begin{exercise}{1.7.17}
Let $G$ be a group and let $G$ act on itself by left conjugation, so each $g\in G$ maps $G$ to $G$ by
\[
    x\mapsto gxg^{-1}.
\]
For fixed $g\in G$, prove that conjugation by $g$ is an isomorphism from $G$ onto itself, i.e., an automorphism of $G$ (cf.\ Exercise 1.6.20). Deduce that $x$ and $gxg^{-1}$ have the same order for all $x$ in $G$ and that for any subset $A$ of $G$,
\[
    |A|=|gAg^{-1}|,
\]
where $gAg^{-1}=\{gag^{-1}\mid a\in A\}$.
\end{exercise}

\begin{solution}
For fixed $g\in G$, define $\varphi_g(x)=gxg^{-1}$. Then
\[
    \varphi_g(xy)=gxyg^{-1}=(gxg^{-1})(gyg^{-1})
    =\varphi_g(x)\varphi_g(y),
\]
so $\varphi_g$ is a homomorphism. Conjugation by $g^{-1}$ is its inverse, since
\[
    \varphi_{g^{-1}}(\varphi_g(x))=x.
\]
Thus $\varphi_g$ is an automorphism. Isomorphisms preserve element orders, so
\[
    |x|=|gxg^{-1}|.
\]
Moreover, the restriction of $\varphi_g$ gives a bijection from $A$ onto $gAg^{-1}$, hence
\[
    |A|=|gAg^{-1}|.
\]
\end{solution}

\begin{exercise}{2.1.2}
In each of (a)--(e), prove that the specified subset is not a subgroup of the given group.
\begin{exercisesubparts}
    \item The set of $2$-cycles in $S_n$ for $n\ge3$.
    \item The set of reflections in $D_n$ for $n\ge3$.
    \item For a composite integer $n>1$ and a group $G$ containing an element of order $n$, the set
    \[
        \{x\in G\mid |x|=n\}\cup\{1\}.
    \]
    \item The set of positive and negative odd integers in $\Z$, together with $0$, under addition.
    \item The set of real numbers whose square is a rational number, under addition.
\end{exercisesubparts}
\end{exercise}

\begin{notation}
Following MATH000, $D_n$ denotes the dihedral group of symmetries of a regular $n$-gon, so $|D_n|=2n$. Dummit--Foote denotes the same group by $D_{2n}$.
\end{notation}

\begin{solution}
\begin{exercisesubparts}
    \item The identity permutation is not a $2$-cycle, so the set does not contain the identity.

    \item The identity is not a reflection, so the set does not contain the identity.

    \item Let $x\in G$ have order $n$. Since $n$ is composite, write $n=st$ with $1<s,t<n$. Then Exercise 1.1.23 gives
    \[
        |x^s|=t.
    \]
    Thus $x$ belongs to the set but $x^s$ does not. A subgroup containing $x$ must contain all powers of $x$, so the set is not a subgroup.

    \item The element $1$ belongs to the set, but
    \[
        1+1=2
    \]
    does not. Hence the set is not closed under addition.

    \item Both $1$ and $\sqrt2$ belong to the set, but
    \[
        (1+\sqrt2)^2=3+2\sqrt2\notin\Q.
    \]
    Hence $1+\sqrt2$ does not belong to the set, so it is not closed under addition.
\end{exercisesubparts}
\end{solution}

\begin{exercise}{2.1.3}
Show that the following subsets of the dihedral group $D_4$ are subgroups:
\[
    \{1,r^2,s,sr^2\},
    \qquad
    \{1,r^2,sr,sr^3\}.
\]
\end{exercise}

\begin{solution}
Let
\[
    H=\{r^{2i},sr^{2i}\mid i=0,1\}.
\]
For $i,j\in\{0,1\}$, using $r^4=s^2=1$ and $r^is=sr^{-i}$,
\[
\begin{aligned}
    r^{2i}r^{2j}&=r^{2(i+j)},&
    r^{2i}(sr^{2j})&=sr^{2(j-i)},\\
    (sr^{2i})r^{2j}&=sr^{2(i+j)},&
    (sr^{2i})(sr^{2j})&=r^{2(j-i)}.
\end{aligned}
\]
Thus $H$ is closed under multiplication. Since $H$ is finite and nonempty, $H\le D_4$.

Similarly, let
\[
    K=\{r^{2i},sr^{2i+1}\mid i=0,1\}.
\]
Then
\[
\begin{aligned}
    r^{2i}r^{2j}&=r^{2(i+j)},&
    r^{2i}(sr^{2j+1})&=sr^{2(j-i)+1},\\
    (sr^{2i+1})r^{2j}&=sr^{2(i+j)+1},&
    (sr^{2i+1})(sr^{2j+1})&=r^{2(j-i)}.
\end{aligned}
\]
Hence $K$ is closed under multiplication, and since $K$ is finite and nonempty, $K\le D_4$.
\end{solution}

\begin{exercise}{2.1.6}
Let $G$ be an abelian group. Prove that
\[
    \{g\in G\mid |g|<\infty\}
\]
is a subgroup of $G$, called the \emph{torsion subgroup} of $G$. Give an explicit example where this set is not a subgroup when $G$ is nonabelian.
\end{exercise}

\begin{solution}
Let
\[
    H=\{g\in G\mid |g|<\infty\}.
\]
Clearly $1\in H$. If $x,y\in H$ with $|x|=m$ and $|y|=n$, then, since $G$ is abelian,
\[
    (xy^{-1})^{mn}=x^{mn}y^{-mn}=1.
\]
Thus $xy^{-1}\in H$, so the Subgroup Criterion gives $H\le G$.

For a nonabelian example, in $\GL_2(\R)$ let
\[
    \mathbf A=\begin{pmatrix}0&1\\1&0\end{pmatrix},
    \qquad
    \mathbf B=\begin{pmatrix}0&2\\\frac12&0\end{pmatrix}.
\]
Then $\mathbf A^2=\mathbf B^2=\mathbf I$, but
\[
    \mathbf A\mathbf B=\begin{pmatrix}\frac12&0\\0&2\end{pmatrix},
    \qquad
    (\mathbf A\mathbf B)^k
    =\begin{pmatrix}2^{-k}&0\\0&2^k\end{pmatrix}\neq\mathbf I
\]
for every $k>0$, so $\mathbf A\mathbf B$ has infinite order. Hence the finite-order elements need not form a subgroup in a nonabelian group.
\end{solution}

\begin{exercise}{2.1.13}
Let $H$ be a subgroup of the additive group of rational numbers with the property that $1/x\in H$ for every nonzero element $x$ of $H$. Prove that
\[
    H=\{0\}\quad\text{or}\quad H=\Q.
\]
\end{exercise}

\begin{solution}
Assume $H\neq\{0\}$ and choose
\[
    0\neq x=\frac ab\in H,
\]
where $b>0$ and $\gcd(|a|,b)=1$. By hypothesis, $1/x=b/a\in H$. Since $H$ is an additive subgroup,
\[
    bx=a\in H,
    \qquad
    a\left(\frac1x\right)=b\in H.
\]
By B\'ezout's identity there exist $m,n\in\Z$ such that $ma+nb=1$, so $1\in H$. Hence every integer belongs to $H$. For any nonzero $q\in\Z$, the hypothesis then gives $1/q\in H$, and therefore
\[
    \frac pq=p\left(\frac1q\right)\in H
\]
for every $p/q\in\Q$. Thus $H=\Q$.
\end{solution}

\begin{exercise}{2.3.1}
Find all subgroups of $C_{45}=\langle x\rangle$, giving a generator for each. Describe the containments between these subgroups.
\end{exercise}

\begin{notation}
Following MATH000, $C_n$ denotes a cyclic group of order $n$. Dummit--Foote denotes the same abstract cyclic group by $Z_n$. The notation $\Z/n\Z$ denotes the residue classes modulo $n$; when used here as a group, it is understood with addition. Residue classes are denoted by $[a]_n$.
\end{notation}

\begin{solution}
The positive divisors of $45$ are
\[
    1,3,5,9,15,45.
\]
A cyclic group has exactly one subgroup of order $d$ for each divisor $d$ of its order, namely $\langle x^{45/d}\rangle$. Thus the subgroups are
\[
\begin{array}{c@{\qquad}c}
\toprule
\text{order} & \text{subgroup}\\
\midrule
1 & \{1\}=\langle x^{45}\rangle\\
3 & \langle x^{15}\rangle\\
5 & \langle x^9\rangle\\
9 & \langle x^5\rangle\\
15 & \langle x^3\rangle\\
45 & \langle x\rangle\\
\bottomrule
\end{array}
\]
If $H_d$ denotes the subgroup of order $d$, then
\[
    H_d\le H_e
    \quad\Longleftrightarrow\quad
    d\mid e.
\]
This completely describes the containments.
\end{solution}

\begin{exercise}{2.3.8}
Let $C_{48}=\langle x\rangle$. For which integers $a$ does the map $\varphi_a$ defined by
\[
    \varphi_a:[1]_{48}\mapsto x^a
\]
extend to an isomorphism from $\Z/48\Z$ onto $C_{48}$?
\end{exercise}

\begin{solution}
For every integer $a$, the relation $48[1]_{48}=[0]_{48}$ is respected because $(x^a)^{48}=1$, so $[1]_{48}\mapsto x^a$ extends to a homomorphism. It is an isomorphism exactly when $x^a$ is a generator of $C_{48}$. For a cyclic group of order $48$,
\[
    |x^a|=\frac{48}{\gcd(a,48)}.
\]
Hence
\[
    \varphi_a\text{ is an isomorphism}
    \quad\Longleftrightarrow\quad
    \gcd(a,48)=1.
\]
Equivalently,
\[
    a\equiv1,5,7,11,13,17,19,23,25,29,31,35,37,41,43,47\pmod{48}.
\]
\end{solution}

\begin{exercise}{2.3.11}
Find all cyclic subgroups of $D_4$. Find a proper subgroup of $D_4$ which is not cyclic.
\end{exercise}

\begin{solution}
The elements of $D_4$ are
\[
    1,r,r^2,r^3,s,sr,sr^2,sr^3,
\]
where $|r|=4$, $|r^2|=2$, and every reflection has order $2$. Hence the distinct cyclic subgroups are
\[
    \{1\},
    \qquad
    \langle r\rangle,
    \qquad
    \langle r^2\rangle,
    \qquad
    \langle s\rangle,
    \qquad
    \langle sr\rangle,
    \qquad
    \langle sr^2\rangle,
    \qquad
    \langle sr^3\rangle.
\]
By Exercise 2.1.3,
\[
    H=\{1,r^2,s,sr^2\}
\]
is a proper subgroup of $D_4$. Every nonidentity element of $H$ has order $2$, so no element generates $H$. Thus $H$ is not cyclic.
\end{solution}

\begin{exercise}{2.3.12}
Prove that the following groups are not cyclic:
\begin{exercisesubparts}
    \item $(\Z/2\Z)\times(\Z/2\Z)$,
    \item $(\Z/2\Z)\times\Z$,
    \item $\Z\times\Z$.
\end{exercisesubparts}
\end{exercise}

\begin{solution}
\begin{exercisesubparts}
    \item Every nonidentity element of $(\Z/2\Z)\times(\Z/2\Z)$ has order $2$. A cyclic group of order $4$ would have an element of order $4$, so this group is not cyclic.

    \item Suppose $(\Z/2\Z)\times\Z=\langle(a,m)\rangle$. Since $(0,1)$ must be generated, there is $k\in\Z$ such that
    \[
        k(a,m)=(0,1),
    \]
    so $km=1$. Hence $m=\pm1$. But then the only multiple of $(a,m)$ with second coordinate $0$ is the zero multiple, so $(1,0)$ cannot be generated. This is a contradiction.

    \item Suppose $\Z\times\Z=\langle(a,b)\rangle$. Since $(1,0)$ must be generated, there is $k\in\Z$ such that
    \[
        k(a,b)=(1,0).
    \]
    Thus $ka=1$ and $kb=0$, so $k=\pm1$ and $b=0$. Then every multiple of $(a,b)$ has second coordinate $0$, so $(0,1)$ cannot be generated. Hence $\Z\times\Z$ is not cyclic.
\end{exercisesubparts}
\end{solution}

\begin{exercise}{2.3.19}
Show that if $H$ is any group and $h$ is an element of $H$, then there is a unique homomorphism from $\Z$ to $H$ such that
\[
    1\mapsto h.
\]
\end{exercise}

\begin{solution}
Define
\[
    \varphi\colon\Z\to H,
    \qquad
    \varphi(n)=h^n.
\]
Then
\[
    \varphi(m+n)=h^{m+n}=h^mh^n=\varphi(m)\varphi(n),
\]
so $\varphi$ is a homomorphism and $\varphi(1)=h$.

If $\psi\colon\Z\to H$ is any homomorphism with $\psi(1)=h$, then for every $n\in\Z$,
\[
    \psi(n)=\psi(1)^n=h^n=\varphi(n).
\]
Thus $\varphi$ is unique.
\end{solution}

\begin{exercise}{2.3.21}
Let $p$ be an odd prime and let $n$ be a positive integer. Use the Binomial Theorem to show that
\[
    (1+p)^{p^{n-1}}\equiv1\pmod{p^n}
\]
but
\[
    (1+p)^{p^{n-2}}\not\equiv1\pmod{p^n}.
\]
Deduce that $1+p$ is an element of order $p^{n-1}$ in the multiplicative group $(\Z/p^n\Z)^\times$.

\textit{Clarification.} The second congruence is understood for $n\ge2$. For $n=1$, the final order statement is immediate.
\end{exercise}

\begin{solution}
For $m\ge0$, we claim that
\[
    (1+p)^{p^m}\equiv1+p^{m+1}\pmod{p^{m+2}}.
\]
The case $m=0$ is immediate. Suppose it holds for $m$, so
\[
    (1+p)^{p^m}=1+p^{m+1}+cp^{m+2}=1+u
\]
for some $c\in\Z$. By the Binomial Theorem,
\[
    (1+u)^p
    =1+pu+\sum_{j=2}^{p-1}\binom pj u^j+u^p.
\]
Since $p$ is prime, $p\mid\binom pj$ for $1\le j\le p-1$, and all terms after $pu$ are divisible by $p^{m+3}$. Hence
\[
    (1+p)^{p^{m+1}}\equiv1+p^{m+2}\pmod{p^{m+3}},
\]
proving the claim.

Taking $m=n-1$ gives
\[
    (1+p)^{p^{n-1}}\equiv1\pmod{p^n}.
\]
For $n\ge2$, taking $m=n-2$ gives
\[
    (1+p)^{p^{n-2}}\equiv1+p^{n-1}\not\equiv1\pmod{p^n}.
\]
Thus, for $n\ge2$, the order of $1+p$ divides $p^{n-1}$ but does not divide $p^{n-2}$, so it is $p^{n-1}$. For $n=1$, $1+p\equiv1\pmod p$, whose order is $1=p^0$.
\end{solution}

\begin{exercise}{2.3.24}
Let $G$ be a finite group and let $x\in G$.
\begin{exercisesubparts}
    \item Prove that if $g\in N_G(\langle x\rangle)$, then
    \[
        gxg^{-1}=x^a
    \]
    for some $a\in\Z$.

    \item Prove conversely that if $gxg^{-1}=x^a$ for some $a\in\Z$, then $g\in N_G(\langle x\rangle)$.

    \textit{Hint.} First show that $gx^kg^{-1}=(gxg^{-1})^k=x^{ak}$ for any integer $k$, so that $g\langle x\rangle g^{-1}\le\langle x\rangle$. If $x$ has order $n$, show that the elements $gx^ig^{-1}$, $i=0,1,\ldots,n-1$, are distinct, so that
    \[
        |g\langle x\rangle g^{-1}|=|\langle x\rangle|=n,
    \]
    and conclude that $g\langle x\rangle g^{-1}=\langle x\rangle$.
\end{exercisesubparts}
\end{exercise}

\begin{solution}
\begin{exercisesubparts}
    \item If $g\in N_G(\langle x\rangle)$, then
    \[
        g\langle x\rangle g^{-1}=\langle x\rangle.
    \]
    Since $x\in\langle x\rangle$, it follows that $gxg^{-1}\in\langle x\rangle$, so
    \[
        gxg^{-1}=x^a
    \]
    for some $a\in\Z$.

    \item Suppose $gxg^{-1}=x^a$. For every $k\in\Z$,
    \[
        gx^kg^{-1}=(gxg^{-1})^k=x^{ak}\in\langle x\rangle,
    \]
    hence
    \[
        g\langle x\rangle g^{-1}\le\langle x\rangle.
    \]
    Let $|x|=n$. If
    \[
        gx^ig^{-1}=gx^jg^{-1}
    \]
    for $0\le i,j<n$, cancellation gives $x^i=x^j$, hence $i=j$. Thus
    \[
        |g\langle x\rangle g^{-1}|=n=|\langle x\rangle|.
    \]
    The subgroup $g\langle x\rangle g^{-1}$ is contained in $\langle x\rangle$ and has the same finite order, so the inclusion is an equality. Hence
    \[
        g\langle x\rangle g^{-1}=\langle x\rangle,
    \]
    so $g\in N_G(\langle x\rangle)$.
\end{exercisesubparts}
\end{solution}

\begin{exercise}{2.3.26}
Let $C_n$ be a cyclic group of order $n$, and for each integer $a$ let
\[
    \sigma_a\colon C_n\to C_n,
    \qquad
    \sigma_a(x)=x^a
\]
for all $x\in C_n$.
\begin{exercisesubparts}
    \item Prove that $\sigma_a$ is an automorphism of $C_n$ if and only if $a$ and $n$ are relatively prime (cf.\ Exercise 1.6.20).
    \item Prove that $\sigma_a=\sigma_b$ if and only if $a\equiv b\pmod n$.
    \item Prove that every automorphism of $C_n$ is equal to $\sigma_a$ for some integer $a$.
    \item Prove that $\sigma_a\circ\sigma_b=\sigma_{ab}$. Deduce that the map
    \[
        [a]_n\mapsto\sigma_a
    \]
    is an isomorphism of $(\Z/n\Z)^\times$ onto $\Aut(C_n)$. Thus $\Aut(C_n)$ is an abelian group of order $\varphi(n)$.
\end{exercisesubparts}
\end{exercise}

\begin{solution}
Let $C_n=\langle z\rangle$.
\begin{exercisesubparts}
    \item Since $C_n$ is abelian,
    \[
        \sigma_a(xy)=(xy)^a=x^ay^a=\sigma_a(x)\sigma_a(y),
    \]
    so $\sigma_a$ is a homomorphism. Its image is $\langle z^a\rangle$. Hence $\sigma_a$ is an automorphism if and only if $z^a$ generates $C_n$, which holds if and only if
    \[
        \gcd(a,n)=1.
    \]

    \item If $\sigma_a=\sigma_b$, then evaluating at $z$ gives $z^a=z^b$, so
    \[
        z^{a-b}=1.
    \]
    Since $|z|=n$, this is equivalent to $n\mid(a-b)$, i.e., $a\equiv b\pmod n$. The converse follows immediately from $x^n=1$ for every $x\in C_n$.

    \item Let $\sigma\in\Aut(C_n)$. Since $z$ has order $n$ and automorphisms preserve element orders, $\sigma(z)$ is a generator of $C_n$. Thus
    \[
        \sigma(z)=z^a
    \]
    for some integer $a$. For every $k\in\Z$,
    \[
        \sigma(z^k)=\sigma(z)^k=z^{ak}=\sigma_a(z^k),
    \]
    so $\sigma=\sigma_a$.

    \item For every $x\in C_n$,
    \[
        (\sigma_a\circ\sigma_b)(x)=\sigma_a(x^b)=x^{ab}=\sigma_{ab}(x).
    \]
    Hence $\sigma_a\circ\sigma_b=\sigma_{ab}$. By parts (a)--(c), the map
    \[
        (\Z/n\Z)^\times\to\Aut(C_n),
        \qquad
        [a]_n\mapsto\sigma_a,
    \]
    is well defined, bijective, and a homomorphism. Therefore
    \[
        (\Z/n\Z)^\times\cong\Aut(C_n).
    \]
    Consequently $\Aut(C_n)$ is abelian and has order $\varphi(n)$.
\end{exercisesubparts}
\end{solution}

\end{document}
